\section{Kinetics and Chemical Equilibrium}

\subsection{Ctrl+F keywords: rate law, reaction order, first order, second order, half-life, Arrhenius, catalyst, equilibrium, Le Chatelier, \(K_p\)}
Use for concentration-time data, identifying order, half-life, activation energy, catalyst statements, equilibrium expressions, \(Q\) versus \(K\), and gas equilibria.

\subsection{Symbols and notation}
\referencetable{tab:kinetics:symbols}{Kinetics and equilibrium symbols}
\begin{tabular}{@{}>{$}l<{$} p{10.8cm}@{}}
\toprule
\text{Symbol} & \text{Meaning and usage}\\
\midrule
r,\ k,\ t,\ [\ce{A}]_t & Reaction rate, rate constant, time, and concentration of \(\ce{A}\) at time \(t\).\\
\dfrac{d[\ce{A}]}{dt} & Change in concentration of species \(\ce{A}\) per unit time.\\
a,b,c,d & Stoichiometric coefficients in the example reaction.\\
m,n & Experimentally determined orders in the rate law; not amounts of substance here.\\
t_{1/2},\ E_{\mathrm a},\ A,\ R,\ T & Half-life, activation energy, frequency factor, gas constant, and temperature.\\
{\mathrm{cat}},\ {\mathrm{uncat}} & Subscripts for catalyzed and uncatalyzed pathways.\\
Q,\ K,\ a_{\ce{A}} & Reaction quotient, equilibrium constant, and thermodynamic activity.\\
\alpha,\beta,\gamma,\delta & Stoichiometric coefficients in the general equilibrium reaction.\\
K_p,\ K_c,\ p_{\ce{A}},\ \Delta n_{\mathrm{gas}} & Pressure/concentration equilibrium constants, partial pressure, and gas coefficient change.\\
\bottomrule
\end{tabular}

\subsection{Kinetics formulas}

\subsubsection{Reaction rate from stoichiometry}
\begin{equation}
a\ce{A}+b\ce{B}\rightarrow c\ce{C}+d\ce{D}
\label{eq:kinetics:reaction-stoichiometry}
\end{equation}
\begin{equation}
r=-\frac{1}{a}\frac{d[\ce{A}]}{dt}
=-\frac{1}{b}\frac{d[\ce{B}]}{dt}
=\frac{1}{c}\frac{d[\ce{C}]}{dt}
=\frac{1}{d}\frac{d[\ce{D}]}{dt}
\label{eq:kinetics:stoichiometric-rate}
\end{equation}
Divide by coefficients so every species reports the same reaction rate.

\subsubsection{Rate law and overall order}
\begin{equation}
r=k[\ce{A}]^m[\ce{B}]^n, \qquad \text{overall order}=m+n
\label{eq:kinetics:rate-law-order}
\end{equation}
Orders are found experimentally unless the problem states an elementary step.

\subsubsection{Initial-rate comparison}
\begin{equation}
\frac{r_2}{r_1}=
\left(\frac{[\ce{A}]_2}{[\ce{A}]_1}\right)^m
\left(\frac{[\ce{B}]_2}{[\ce{B}]_1}\right)^n
\label{eq:kinetics:initial-rate-ratio}
\end{equation}
Compare experiments in which all but one reactant concentration are unchanged to isolate one order at a time.

\subsubsection{Integrated rate laws and linearity test}
\referencetable{tab:kinetics:integrated-rate-laws}{Transformed concentration-time tests for reaction order}
{\small
\begin{tabular}{@{}c l l l@{}}
\toprule
\text{Order} & \text{Transform to plot against \(t\)} & \text{Linear relation} & \text{Slope gives \(k\)}\\
\midrule
0 & \([\ce{A}]_t\) & \([\ce{A}]_t=-kt+[\ce{A}]_0\) & \(k=-\text{slope}\)\\
1 & \(\ln[\ce{A}]_t\) & \(\ln[\ce{A}]_t=-kt+\ln[\ce{A}]_0\) & \(k=-\text{slope}\)\\
2 & \(1/[\ce{A}]_t\) & \(1/[\ce{A}]_t=kt+1/[\ce{A}]_0\) & \(k=\text{slope}\)\\
\bottomrule
\end{tabular}
}
The transformed data set that forms a straight line identifies the order; its fitted slope gives \(k\).

\subsubsection{Half-life}
\begin{equation}
t_{1/2}^{(0)}=\frac{[\ce{A}]_0}{2k}, \qquad
t_{1/2}^{(1)}=\frac{\ln 2}{k}, \qquad
t_{1/2}^{(2)}=\frac{1}{k[\ce{A}]_0}
\label{eq:kinetics:half-life}
\end{equation}
Only first-order half-life is independent of initial concentration.

\subsubsection{Arrhenius equation}
\begin{equation}
k=Ae^{-E_{\mathrm a}/RT}, \qquad
\ln k=-\frac{E_{\mathrm a}}{R}\frac{1}{T}+\ln A
\label{eq:kinetics:arrhenius}
\end{equation}
A plot of \(\ln k\) versus \(1/T\) has slope \(-E_{\mathrm a}/R\).

\subsubsection{Two-temperature Arrhenius relation}
\begin{equation}
\ln\left(\frac{k_2}{k_1}\right)
=-\frac{E_{\mathrm a}}{R}\left(\frac{1}{T_2}-\frac{1}{T_1}\right)
\label{eq:kinetics:arrhenius-two-temperature}
\end{equation}
Use kelvin. If \(T_2>T_1\), then normally \(k_2>k_1\).

\subsubsection{Catalyst effect}
\begin{equation}
E_{\mathrm a,cat}<E_{\mathrm a,uncat}, \qquad
K_{\mathrm{eq,cat}}=K_{\mathrm{eq,uncat}}, \qquad
\Delta G^\circ_{\mathrm{cat}}=\Delta G^\circ_{\mathrm{uncat}}
\label{eq:kinetics:catalyst-effect}
\end{equation}
A catalyst speeds forward and reverse reactions by lowering activation barriers; it does not move equilibrium.

\subsection{Exam recipe: reaction order from concentration-time data}
\textbf{Use when:} a table or graph gives \([\ce{A}]\) versus time and the task asks for order, \(k\), or a concentration prediction.

\textbf{Given / Find:} given \(t\) and \([\ce{A}]_t\); find zero-, first-, or second-order behavior and the corresponding \(k\).

\textbf{Required references:} Table~\ref{tab:kinetics:integrated-rate-laws} and equation~\eqref{eq:kinetics:half-life} when a half-life is requested.

\pythontool{\texttt{kinetic\_linear\_fit(...)} for hver af orden 0, 1 og 2, saa \(R^2\), haeldning og skaering sammenlignes ensartet}

\textbf{Steps:}
\begin{enumerate}
    \item Construct all three transformed data sets by calculating \([\ce{A}]_t\), \(\ln[\ce{A}]_t\), and \(1/[\ce{A}]_t\) for each listed time.
    \item Determine the order by plotting each transformed data set against \(t\), or by comparing successive slopes, and choosing the one with approximately constant slope (best linearity).
    \item Obtain \(k\) from the linear fit slope by using \(k=-\text{slope}\) for the \([\ce{A}]_t\) or \(\ln[\ce{A}]_t\) line and \(k=\text{slope}\) for the \(1/[\ce{A}]_t\) line.
    \item Predict a missing concentration or time by inserting the fitted \(k\), the initial value, and the requested variable into the selected linear relation in Table~\ref{tab:kinetics:integrated-rate-laws}.
\end{enumerate}

\textbf{Checks:} verify that \(k>0\), that only one transformation is clearly linear within data precision, and that first-order data show approximately equal fractional loss over equal time intervals.

\textbf{Common traps:} selecting an order from the raw decreasing curve without testing transformations; using the negative slope itself as \(k\) for order zero or one; taking \(\ln\) of a value with inconsistent concentration units.

\subsection{Exam recipe: reaction order from initial-rate experiments}
\textbf{Use when:} a data table gives initial concentrations and initial rates for multiple experimental runs.

\textbf{Given / Find:} Given initial-rate runs; find the order in each reactant, overall order, or the rate constant.

\textbf{Required references:} Equations~\eqref{eq:kinetics:rate-law-order} and \eqref{eq:kinetics:initial-rate-ratio}.

\textbf{Steps:}
\begin{enumerate}
    \item Choose two runs in which every concentration except one reactant is identical; this makes all unchanged concentration ratios equal to one in Equation~\eqref{eq:kinetics:initial-rate-ratio}.
    \item Calculate the remaining concentration ratio and rate ratio, then solve \(\text{rate ratio}=(\text{concentration ratio})^{\text{order}}\); for ratios not obvious by inspection, use \(\text{order}=\log(r_2/r_1)/\log(c_2/c_1)\).
    \item Repeat step 1 for each reactant whose order is required and add the exponents for overall order using Equation~\eqref{eq:kinetics:rate-law-order}.
    \item If \(k\) is requested, substitute one complete experimental run and the determined exponents into Equation~\eqref{eq:kinetics:rate-law-order}, then isolate \(k\).
\end{enumerate}

\textbf{Checks:} Substitution of a second run should reproduce its rate; unchanged concentrations must not affect the isolated exponent.

\textbf{Common traps:} comparing runs where two concentrations change before both orders are known; treating stoichiometric coefficients as measured orders; reversing numerator and denominator in only one ratio.

\subsection{Exam recipe: Arrhenius relation and catalyst}
\textbf{Use when:} the task gives rate constants at temperatures, asks for \(E_{\mathrm a}\), or asks what a catalyst changes.

\textbf{Given / Find:} given \(k\) and \(T\) values and/or catalyst information; find \(E_{\mathrm a}\), a new \(k\), or the correct catalyst conclusion.

\textbf{Required references:} equations~\eqref{eq:kinetics:arrhenius}, \eqref{eq:kinetics:arrhenius-two-temperature}, and \eqref{eq:kinetics:catalyst-effect}.

\pythontool{\texttt{arrhenius\_ratio(...)} naar opgaven spoerger efter \(k_2/k_1\)}

\textbf{Steps:}
\begin{enumerate}
    \item Prepare temperatures by converting each temperature to kelvin before calculating \(1/T\) or using a temperature ratio relation.
    \item Find \(E_{\mathrm a}\) from two measurements by substituting \(k_1,k_2,T_1,T_2\) into equation~\eqref{eq:kinetics:arrhenius-two-temperature} and isolating \(E_{\mathrm a}\), or from several measurements by fitting \(\ln k\) versus \(1/T\) and using \(E_{\mathrm a}=-R(\text{slope})\).
    \item Find an unknown rate constant by rearranging equation~\eqref{eq:kinetics:arrhenius-two-temperature} after \(E_{\mathrm a}\) is known.
    \item Evaluate a catalyst statement by applying equation~\eqref{eq:kinetics:catalyst-effect}: lower activation energy and faster approach to equilibrium, but unchanged \(K_{\mathrm{eq}}\) and \(\Delta G^\circ\).
\end{enumerate}

\textbf{Checks:} for a positive \(E_{\mathrm a}\), raising \(T\) must raise \(k\); activation energy should have energy-per-mole units.

\textbf{Common traps:} using degrees Celsius in reciprocal temperature; losing the minus sign in the Arrhenius slope; claiming that a catalyst changes equilibrium composition.

\subsection{Equilibrium formulas}

\subsubsection{Reaction quotient and equilibrium constant}
\begin{equation}
\alpha\ce{A}+\beta\ce{B}\rightleftharpoons\gamma\ce{C}+\delta\ce{D}
\label{eq:kinetics:equilibrium-reaction}
\end{equation}
\begin{equation}
Q=\frac{a_{\ce{C}}^\gamma a_{\ce{D}}^\delta}{a_{\ce{A}}^\alpha a_{\ce{B}}^\beta},
\qquad Q_{\mathrm{eq}}=K
\label{eq:kinetics:reaction-quotient}
\end{equation}
Pure solids, pure liquids, and the solvent are omitted from the expression.

\subsubsection{Direction from \(Q\) and \(K\)}
\begin{equation}
Q<K:\ \text{forward shift}, \qquad
Q=K:\ \text{equilibrium}, \qquad
Q>K:\ \text{reverse shift}
\label{eq:kinetics:q-versus-k}
\end{equation}
The system shifts so that \(Q\) approaches \(K\).

\subsubsection{Gas equilibrium \(K_p\)}
\begin{equation}
K_p=\frac{(p_{\ce{C}})^\gamma(p_{\ce{D}})^\delta}{(p_{\ce{A}})^\alpha(p_{\ce{B}})^\beta}
\label{eq:kinetics:kp-expression}
\end{equation}
Use partial pressures with stoichiometric powers and omit pure condensed phases.

\subsubsection{Relation between \(K_p\) and \(K_c\)}
\begin{equation}
K_p=K_c(RT)^{\Delta n_{\mathrm{gas}}},
\qquad
\Delta n_{\mathrm{gas}}=\sum\nu_{\mathrm{gas,products}}-\sum\nu_{\mathrm{gas,reactants}}
\label{eq:kinetics:kp-kc-relation}
\end{equation}

\subsubsection{Le Chatelier shift rules}
\referencetable{tab:kinetics:shift-rules}{Quick equilibrium shift decisions}
\begin{tabular}{@{}l l@{}}
\toprule
\text{Disturbance} & \text{Shift decision}\\
\midrule
Add reactant or remove product & Toward products\\
Remove reactant or add product & Toward reactants\\
Decrease volume for gas system & Toward fewer moles of gas\\
Raise \(T\) for an endothermic forward reaction & Toward products\\
Raise \(T\) for an exothermic forward reaction & Toward reactants\\
Add catalyst & No equilibrium shift\\
\bottomrule
\end{tabular}

\subsection{Exam recipe: equilibrium expression and shift}
\textbf{Use when:} the task asks for \(K\), \(Q\), \(K_p\), reaction direction, or a Le Chatelier/catalyst statement.

\textbf{Given / Find:} given a balanced equilibrium and activities, concentrations, partial pressures, or a disturbance; find the equilibrium expression, \(Q\)-based direction, or qualitative shift.

\textbf{Required references:} equations~\eqref{eq:kinetics:reaction-quotient}, \eqref{eq:kinetics:q-versus-k}, \eqref{eq:kinetics:kp-expression}, and \eqref{eq:kinetics:kp-kc-relation}; Table~\ref{tab:kinetics:shift-rules}.

\textbf{Steps:}
\begin{enumerate}
    \item Write the expression by placing product activities over reactant activities and using balanced stoichiometric coefficients as exponents, while omitting pure solids and pure liquids.
    \item Choose the quantity type by using partial pressures in equation~\eqref{eq:kinetics:kp-expression} for \(K_p\), concentrations for \(K_c\), or converting with equation~\eqref{eq:kinetics:kp-kc-relation} when required.
    \item Determine the spontaneous direction from measured values by calculating \(Q\) with the same expression as \(K\) and comparing it with \(K\) using equation~\eqref{eq:kinetics:q-versus-k}.
    \item Determine the response to an imposed change by locating the disturbance in Table~\ref{tab:kinetics:shift-rules}; for temperature changes, treat heat as a reactant for an endothermic forward reaction and as a product for an exothermic one; for a catalyst, state faster equilibration rather than a shifted position.
\end{enumerate}

\textbf{Checks:} exponents must match the balanced equation; \(Q=K\) means no net shift; a volume change only favors a side if the gaseous mole totals differ.

\textbf{Common traps:} including a pure solid or liquid in \(K\); reversing the \(Q<K\) direction; stating that a catalyst changes \(K\).

\subsubsection{ICE-table extent relation}
\begin{equation}
n_i=n_{i,0}+\nu_i\xi, \qquad [i]=\frac{n_i}{V}
\label{eq:kinetics:ice-extent}
\end{equation}
Use negative signed coefficients \(\nu_i\) for reactants and positive coefficients for products in the forward direction.

\subsection{Exam recipe: equilibrium composition from \(K\) and initial amounts}
\textbf{Use when:} initial amounts or concentrations and an equilibrium constant are given, and the equilibrium amount or concentration is requested.

\textbf{Given / Find:} Given a balanced equilibrium, initial composition, volume, and \(K\) or thermodynamic data from which \(K\) can be obtained; find the equilibrium composition.

\textbf{Required references:} Equations~\eqref{eq:kinetics:reaction-quotient}, \eqref{eq:kinetics:ice-extent}, and \eqref{eq:thermo:equilibrium-constant} when \(K\) must first be calculated.

\pythontool{\texttt{equilibrium\_constant(...)} naar \(K\) foerst skal findes fra \(\Delta G^\circ\); selve ICE-polynomiet skal opstilles fra reaktionen foer numerisk loesning}

\textbf{Steps:}
\begin{enumerate}
    \item Balance the equilibrium and write the \(K\) expression from Equation~\eqref{eq:kinetics:reaction-quotient}; omit pure solids and liquids.
    \item Obtain \(K\): use the supplied value directly or calculate it from \(\Delta G^\circ\) with Equation~\eqref{eq:thermo:equilibrium-constant}.
    \item Define one forward extent \(\xi\) and write every equilibrium amount as \(n_{i,0}+\nu_i\xi\) using Equation~\eqref{eq:kinetics:ice-extent}; divide by common volume when a concentration expression is needed.
    \item Substitute those expressions into \(K\), solve for \(\xi\), and reject any mathematical root that makes any physical amount negative.
    \item Insert the accepted \(\xi\) into each requested species expression and report amount or concentration with units.
\end{enumerate}

\textbf{Checks:} Equilibrium amounts are nonnegative and give the stated \(K\) when substituted back; the sign of \(\xi\) agrees with the initial \(Q\)-versus-\(K\) direction.

\textbf{Common traps:} using initial concentrations directly as equilibrium values; omitting stoichiometric multipliers in the ICE row; retaining a nonphysical algebraic root.
